\section{Task 1: Frequency-based Baseline}

\begin{questionbox}
\textbf{Q1.} How many unique API calls does the training set contain? How many the test set?
\end{questionbox}

The training set contains 258 unique API calls, while the test set contains 232 unique API calls.
The train partition has 16,325 samples and the test partition has 6,505 samples. Both partitions
are strongly imbalanced: malware represents about 96.26\% of the samples in both train and test.

\begin{questionbox}
\textbf{Q2.} Are there any API calls that appear only in the test set (but not in the training set)?
If yes, how many? And which ones are they?
\end{questionbox}

Yes. There are 3 API calls that appear only in the test set:
\texttt{ControlService}, \texttt{NtDeleteKey}, and \texttt{WSASocketA}. Each of them appears once
in the test set. The intersection between train and test vocabularies contains 229 API calls.

\begin{questionbox}
\textbf{Q3.} Can you use the test vocabulary to build the new test dataframe?
If not, how do you handle API calls in the test set that do not exist in the training vocabulary?
\end{questionbox}

No. The test vocabulary should not be used to define the feature space, because it would introduce
information from the test set into preprocessing. We build both the train and test frequency matrices
using only the training vocabulary. Therefore, API calls that appear only in the test set are treated
as out-of-vocabulary terms and ignored, i.e. they do not get a dedicated feature column.

\begin{questionbox}
\textbf{Q4.} One issue of this frequency-based approach is that it creates sparse vectors.
How many non-zero elements per row do you have on average in the training set?
How many in the test set?
What is the ratio with respect to the number of elements per row?
\end{questionbox}

The frequency vectors have 258 features, one for each API call in the training vocabulary.
The average number of non-zero elements per row is 21.9471 in the training set and 24.2782 in
the test set. This corresponds to non-zero ratios of 0.0851 and 0.0941 respectively.
Equivalently, about 91.49\% of the train vectors and 90.59\% of the test vectors are zeros.

\begin{table}[h!]
\centering
\begin{tabular}{lcccc}
\hline
\textbf{Split} & \textbf{Features} & \textbf{Avg. non-zero} & \textbf{Non-zero ratio} & \textbf{Zero ratio} \\
\hline
Train & 258 & 21.9471 & 0.0851 & 0.9149 \\
Test  & 258 & 24.2782 & 0.0941 & 0.9059 \\
\hline
\end{tabular}
\caption{Sparsity of the frequency-based representation.}
\label{tab:task1_sparsity}
\end{table}

\begin{questionbox}
\textbf{Q5.} The original API sequences were ordered. Is it still the case now (i.e., in the frequency-based dataframe)? Why?
\end{questionbox}

No. The frequency-based dataframe only stores how many times each API call appears in a sequence.
It does not store positions or transitions, so the original temporal order is lost.

\begin{questionbox}
\textbf{Q6.} Report how you chose the hyperparameters of your classifier, and the final performance on the test set.
\end{questionbox}

We used Logistic Regression as the frequency-based baseline classifier. Hyperparameters were selected
with Optuna using 5-fold stratified cross-validation on the training set, optimizing macro F1. The
search considered the hyperparameter space reported in Table~\ref{tab:task1_optuna_space}.

\begin{table}[h!]
\centering
\begin{tabular}{lll}
\hline
\textbf{Hyperparameter} & \textbf{Search space} & \textbf{Scale/type} \\
\hline
\texttt{C} & $[10^{-3}, 100]$ & Log-uniform \\
\texttt{class\_weight} & \{\texttt{None}, \texttt{balanced}\} & Categorical \\
\texttt{penalty} & \{\texttt{l1}, \texttt{l2}\} & Categorical \\
\hline
\end{tabular}
\caption{Optuna search space for the Logistic Regression baseline.}
\label{tab:task1_optuna_space}
\end{table}

The best configuration found was:
\[
C = 41.3185, \qquad \texttt{class\_weight}=\texttt{None}, \qquad \texttt{penalty}=\texttt{l2}.
\]
The best mean macro F1 during cross-validation was 0.7180.

\begin{table}[h!]
\centering
\begin{tabular}{lcccc}
\hline
\textbf{Class} & \textbf{Precision} & \textbf{Recall} & \textbf{F1-score} & \textbf{Support} \\
\hline
Goodware (0) & 0.6613 & 0.3374 & 0.4469 & 243 \\
Malware (1)  & 0.9748 & 0.9933 & 0.9839 & 6262 \\
\hline
Accuracy     &        &        & 0.9688 & 6505 \\
Macro avg    & 0.8180 & 0.6654 & 0.7154 & 6505 \\
Weighted avg & 0.9631 & 0.9688 & 0.9639 & 6505 \\
\hline
\end{tabular}
\caption{Classification report of the Logistic Regression baseline on the test set.}
\label{tab:task1_classification_report}
\end{table}

\begin{figure}[h!]
\centering
\includegraphics[width=0.62\textwidth]{figures/task-1/confusion_matrix.pdf}
\caption{Confusion matrix of the Logistic Regression baseline on the test set.}
\label{fig:task1_confusion_matrix}
\end{figure}

The global test metrics were: accuracy 0.9688, balanced accuracy 0.6654, macro precision 0.8180,
macro recall 0.6654, and macro F1 0.7154.

\begin{questionbox}
\textbf{Q7.} Is the final performance good, even ignoring the order of API calls and handling very sparse vectors?
\end{questionbox}

The baseline is strong at detecting malware: malware recall is 0.9933 and malware F1 is 0.9839.
However, the high accuracy is partly explained by the strong class imbalance. The model performs
much worse on goodware, with recall 0.3374 and F1 0.4469, meaning that many goodware samples are
classified as malware. Therefore, the frequency-based Logistic Regression is a useful and competitive
baseline, but it is not fully satisfactory for a real-world detector because its performance on the
minority goodware class is limited.
